% 从仓库根目录构建：
% mkdir -p .build/midterm/briefing-tex
% xelatex -halt-on-error -interaction=nonstopmode -output-directory=.build/midterm/briefing-tex presentations/02-midterm/briefing/第11组_中期组员准备手册.tex
\documentclass[UTF8,11pt,fontset=fandol]{ctexart}
\usepackage[a4paper,margin=23mm,headheight=23pt,headsep=9mm,footskip=12mm]{geometry}
\usepackage{xcolor,enumitem,fancyhdr}
\usepackage{hyperref}
\definecolor{zblue}{HTML}{003F88}
\definecolor{pale}{HTML}{EDF3FA}
\definecolor{ink}{HTML}{233042}
\definecolor{muted}{HTML}{52647C}
\definecolor{ruleblue}{HTML}{B9CCE3}
\hypersetup{colorlinks=true,linkcolor=zblue,urlcolor=zblue,pdftitle={第11组中期组员准备手册},pdfauthor={第11组}}
\pagestyle{fancy}
\fancyhf{}
\lhead{\small\color{zblue}软件工程管理 \quad 第11组}
\rhead{\small\color{muted}中期汇报组内准备}
\rfoot{\small\color{muted}\thepage}
\renewcommand{\headrulewidth}{0.55pt}
\renewcommand{\headrule}{\hbox to\headwidth{\color{ruleblue}\leaders\hrule height \headrulewidth\hfill}}
\setlength{\parindent}{0pt}
\setlength{\parskip}{5pt}
\setlist[itemize]{leftmargin=1.7em,topsep=3pt,itemsep=4pt}
\newcommand{\ptitle}[1]{\vspace*{3pt}{\color{zblue}\LARGE\bfseries #1}\par\vspace{8pt}}
\newcommand{\subtopic}[1]{\par\vspace{5pt}{\color{zblue}\large\bfseries #1}\par\vspace{1pt}}
\newcommand{\captionline}[1]{{\small\color{muted}#1}\par}
\newcommand{\lead}[1]{\noindent\colorbox{pale}{\parbox{\dimexpr\linewidth-2\fboxsep\relax}{\vspace{4pt}\raggedright #1\vspace{4pt}}}\par}
\begin{document}
\color{ink}

\ptitle{中期组员准备手册}
{\large 第11组 \quad 产业智能报告生成系统}\par
\captionline{组长兼主讲人 \quad 程文博}

\lead{\textbf{中期先做好一份光伏简报。} 系统用工信部和国家能源局的公开资料，写2024年中国光伏的“制造与价格”和“发电新增装机”两章。缺数据时标“待核实”。人核对草稿和原文后，再确认导出 PDF。}

\subtopic{老师上次提的三件事}

\begin{itemize}
\item \textbf{参考现有框架。} 我们看了 DeepSeek Harness 和 Pi CLI，分别用本地模拟模型检查同一个资料工具能否被调用、结果能否送回模型。chat-demo 用来参考附件和草稿怎么展示。
\item \textbf{把范围收窄。} 开题想覆盖多类产业报告，中期先用光伏两章检验数据能否回查、缺资料时报告会怎么写。
\item \textbf{按智能体的工作分工。} 定章节、找资料、接工具、核对报告、复跑演示，大家围绕同一份简报交接。
\end{itemize}

\subtopic{现在已经做成什么}

\begin{itemize}
\item 在命令行用 LangGraph 和 GLM 跑通了查资料、写草稿、核对和导出。两份资料齐全时有3项指标可回查。只移除装机来源时，制造数据还在，装机一栏标为“待核实”。
\item 数字由程序从证据中填入，模型辅助写观察。这样读者能找到数字的出处，资料不足也一眼看得出来。
\item 已验证写完第一章后中断，重启时能跳过第一章。后续还用固定样例测试了资料冲突、客户端超时和风电模板复用。风电用于验证复用，中期主讲光伏。
\end{itemize}

\subtopic{先看两份样张，再开始准备}

打开仓库里的“完整资料”和“缺装机资料”两份 PDF，对比“发电新增装机”一栏。写报告、做 PPT 和跑演示都用这组例子，大家对着同一份结果核对。
\captionline{\href{https://github.com/BrockPhlin/industry-report-harness/tree/codex/midterm-prep/presentations/02-midterm}{打开仓库里的中期准备材料和样张}}

\newpage

\ptitle{每个人具体做什么}
\lead{完成后把文件或 PR 链接发给程文博，带一句你改了什么、核对了什么。}

\subtopic{程文博｜统筹和主讲}
收齐报告、PPT 和样张，核对讲稿中的数字、术语和实验结果。组织一次计时彩排，负责正式汇报。组员各自准备负责部分的问答。

\subtopic{田亦洁｜写系统设计报告}
按仓库里的报告框架写和改。把总体架构图、光伏生成流程、输入输出和完整／缺资料两组结果写进报告。交可编辑稿和 PDF。

\subtopic{崔大元｜整理资料、核对数字}
整理一张来源表，核对组件产量588 GW、组件价格同比下降29.7\%和新增装机2.78亿千瓦。每条带上单位、统计期、原文链接和所在段落，把表交给田亦洁和李佳雨。

\subtopic{许子萌｜整理架构图和框架说明}
给田亦洁和李佳雨一张流程图，画清模板、查资料、补查、写草稿、人工确认和导出。再用两段话写明 DeepSeek Harness 怎么接工具、Pi CLI 怎么记录工具调用，附一条可核对的运行记录。

\subtopic{李佳雨｜完善中期 PPT}
按仓库里的内容框架完善 PPT。把老师的三条意见、系统流程、两份 PDF 的同栏对照和分工放进去。用崔大元的来源表核对数字，给程文博一份可编辑的 PPTX，便于一起调整。

\subtopic{郭润航｜复跑项目、准备演示}
本地跑一遍完整资料和缺装机资料，两组都导出 PDF。把运行命令、实际结果和两份样张交给程文博。两份样张也发给李佳雨。彩排时检查离线能否复现，负责计时。

\end{document}
